% ECONOMICS_PROBLEM: {"id": "J62-351", "title": "Causality behind the Great Gatsby curve", "jel_code": "J62", "source_id": "OP-0368"}
% !TeX program = xelatex
\documentclass[11pt,a4paper]{article}
\usepackage[margin=23mm]{geometry}
\usepackage{fontspec}
\setmainfont{Latin Modern Roman}
\usepackage{amsmath,amssymb,mathtools}
\usepackage{xcolor,enumitem,booktabs,longtable,array,xurl,hyperref}
\definecolor{muted}{HTML}{53636C}
\hypersetup{colorlinks=true,urlcolor=blue,linkcolor=blue}
\setlength{\parindent}{0pt}
\setlength{\parskip}{5pt}
\newcommand{\R}{\mathbb R}
\newcommand{\E}{\mathbb E}
\newcommand{\N}{\mathbb N}
\newcommand{\ind}{\mathbf 1}
\newcommand{\Var}{\operatorname{Var}}
\newcommand{\Cov}{\operatorname{Cov}}
\newcommand{\supp}{\operatorname{supp}}
\DeclareMathOperator*{\argmin}{arg\,min}
\DeclareMathOperator*{\argmax}{arg\,max}
\newcommand{\entrymeta}[3]{{\sffamily\small\color{muted}\textbf{\texttt{#1}}\quad|\quad #2\par #3\par}}
\newcommand{\Problemref}[1]{\texttt{#1}}
\newcommand{\JELref}[2]{\texttt{#2} (JEL \texttt{#1})}
\begin{document}
\section*{Causality behind the Great Gatsby curve}
\textbf{Problem ID:} \texttt{J62-351}\quad \textbf{JEL:} \texttt{J62}\par
% BEGIN ECONOMICS STATEMENT
\label{problem:OP-0368}
{\sffamily\small\color{muted}Original subject group: Inequality and economic mobility\par}
\label{definitions:problem:30007650}
\textbf{Audit classification:} Background; proposed extension. Authors, title, 2022 volume and DOI checked. Abstract summarizes multiple models and findings; the exact policy-causal problem is not certified open. Full-text agenda not inspected.
\subsection*{Definitions and precise research target}
\textbf{Complete editorial problem.} Fix a collection of regions, a parental cohort, and a child cohort. Let \(F^P_r\) and \(F^C_r\) be the parent and child income distributions in region \(r\). Define inequality \(I_r\) with a predeclared functional, and define persistence \(\beta_r=\operatorname{Cov}(\log Y^P,\log Y^C)/\operatorname{Var}(\log Y^P)\) for positive lifetime incomes with finite log-income second moments and strictly positive parental log-income variance under every compared policy. The observed association between \(I_r\) and \(\beta_r\) is the Great Gatsby relationship, not a causal estimand. Let \(p\) be a specified intervention changing parental inequality, with stated effects on mean income and public resources. Define \(\tau_r(p)=\beta_r(p)-\beta_r(0)\). Identify these policy effects and determine whether they arise through family investments, skills, social influences, political choices, or aspirations. Different interventions producing the same inequality statistic may have different effects; therefore \(\beta\) is not assumed to be a function of inequality alone. Account for changing marginal distributions, measurement error, and regional composition that can create mechanical relationships. The task is a proposed causal specification informed by a review of multiple mechanisms. The reviewed abstract does not itself certify the catalogue's broad causal formulation as an explicitly unsolved problem.

\textbf{Definitions and admissible scope.} For every regime, lifetime incomes are positive with finite log second moments and positive parental log variance. Persistence \(\beta\) is the coefficient of the population linear projection of child log income on parent log income, not a correlation unless variances are equal. Region inequality \(I\) is computed from a separately fixed population and income concept. Intervention \(p\) specifies who receives resources, fiscal finance, public investment, and timing; changing a Gini is not itself an intervention. Compare \(\beta(p)-\beta(0)\) after specifying whether ranks, logs, migration, and cohorts are held comparable. If \(\log Y^C=a+\rho\log Y^P+\epsilon\), with uncorrelated error, then \(\rho\) can remain fixed when parental inequality changes; an observed curve is therefore insufficient for a causal inequality elasticity. The parent--child linkage rule and paired-family weights are fixed across compared regimes; \(F^P_r,F^C_r\) alone do not determine the joint covariance. Use the same real income concept and ages, and specify \(I\), for example a nonnegative-income Gini. An admissible data-generating class contains joint parent--child laws and potential outcomes under every feasible \(p\), with stated migration, selection, fiscal, and causal-assignment restrictions. Return the identified set of \((I_r(p),\beta_r(p),\tau_r(p))\). A statement relating aggregate region statistics is descriptive unless the separate policy contrasts are identified.
\subsection*{Variants and assumption changes}
\begin{enumerate}
\item Rank mobility variant (editorial): replace log-income elasticity by \(\operatorname{Cov}(R^P,R^C)/\operatorname{Var}(R^P)\), with ranks and tie rules defined against fixed reference distributions. This reduces scale dependence but does not identify causality.
\item Policy-equivalence test (editorial): compare two equal-budget interventions that produce the same parental Gini but different education or network changes. Test whether their persistence effects differ, rejecting a scalar sufficient-statistic interpretation when they do.
\end{enumerate}
\subsection*{References and source audit}
\textbf{Support and access.} Authors, title, 2022 volume and DOI checked. Abstract summarizes multiple models and findings; the exact policy-causal problem is not certified open. Full-text agenda not inspected. \textbf{Locator:} Publisher abstract, paragraphs 1--2.
\begin{enumerate}\raggedright
\item\hypertarget{definitions:source:30007650:1}{} Steven N. Durlauf; Andros Kourtellos; Chih Ming Tan. 2022. \emph{The Great Gatsby Curve}. Abstract, paragraphs 1–2.
\url{https://www.annualreviews.org/content/journals/10.1146/annurev-economics-082321-122703}.
DOI: \url{https://doi.org/10.1146/annurev-economics-082321-122703}.
\end{enumerate}

\subsection*{Common conventions: Source mathematical conventions}

These conventions apply to each entry unless explicitly replaced.

\textbf{Models and identification.} A primitive model \(m\) specifies all probability laws, feasible choices, information, equilibrium and measurement rules used by its target. The admissible class \(\mathcal M\) is fixed independently of outcomes. If \(P_O(m)\) is its observable probability law and \(\tau(m)\) the target, its identified set at observable law \(P\) is
\[
 \mathcal I_\tau(P)=\{\tau(m):m\in\mathcal M,\ P_O(m)=P\}.
\]
Point identification means that this set is a singleton. An empty set signals incompatibility of data and assumptions. Coordinate bounds need not describe the joint set. Bounds are sharp only if every retained target is attainable or arbitrarily approachable by compatible models. Finite bounds require explicit restrictions on unobserved quantities; unrestricted confounding, hidden wealth, extrapolation or arbitrary equilibrium selection can yield uninformative sets.

\textbf{Causal interventions.} A potential outcome \(Y_i(p)\) is the outcome under a complete policy regime \(p\), including timing, financing, information and market spillovers. The target population and units are fixed before treatment. With interference, use the complete assignment vector or a justified exposure mapping; individual no-interference notation alone cannot identify an economy-wide intervention. A difference of mean potential outcomes does not require the cross-world joint distribution. Conditioning on post-treatment migration, survival, enrollment or employment changes the estimand and needs separate justification.

\textbf{Statistical statements.} An estimator, confidence set, forecast or test specifies a sampling law, model class and loss/coverage criterion. Uniform \((1-\alpha)\)-coverage means \(\inf_{m\in\mathcal M}\Pr_m\{\tau(m)\in C_n\}\geq1-\alpha\), or an explicitly asymptotic version. The whole parameter space trivially has coverage, so informativeness requires a stated width, risk or power objective. A forecasting improvement is relative to a named benchmark, information set and evaluation population.

\textbf{Equilibrium and optimization.} An equilibrium comprises feasible plans, optimality, beliefs where needed, budget constraints and all market/resource clearing. The primitives or a stated benchmark define each correspondence \(\mathcal E(p;m)\). Nonemptiness is assumed or proved, never inferred just from compact policy sets. A selected-equilibrium target uses a declared selection rule; a robust target quantifies over all permitted equilibria. Use suprema or infima unless attainment is established. An \(\varepsilon\)-optimal policy is feasible and within \(\varepsilon\) of the finite optimum value. Report an empty feasible set separately.

\textbf{Welfare and accounting.} Normative utility, interpersonal normalization, social weights, numeraire, discounting and the use of public receipts are primitives. Behavioral utility need not coincide with the welfare criterion. Household welfare includes ownership income and rents once; adding profits again to compensating variation generally double-counts them. A monetary equivalent is defined by an explicit transfer and price convention, with monotonicity and range conditions. Average effects, mechanism contributions, transfers and welfare losses are different targets. Infinite sums require convergence and dynamic feasible plans require terminal or no-Ponzi conditions.

\textbf{Source versions.} A working-paper date, revision date and journal publication year are distinguished. A DOI for a journal version is not automatically the DOI of an earlier manuscript. Locators refer to the stated version and printed pages where specified. A metadata-only check does not verify a claim about a full-text conclusion. Every entry's source subsection narrows the provenance claim accordingly.
% END ECONOMICS STATEMENT
\end{document}
